\documentclass[12pt]{article}
\usepackage[utf8]{inputenc}
\usepackage[T1]{fontenc}
\usepackage{times}
\usepackage[margin=2cm]{geometry}
\usepackage{setspace}

\setstretch{1.5}

\begin{document}

\begin{center}
    \textbf{\Large CV}
\end{center}

\section*{Family name, First name(s)}
Nielsen, Finn Årup

\section*{Current position(s)}
2006--Present \quad Researcher/Professor, Department of [Insert Dept], Technical University of Denmark, Denmark

\section*{Previous positions}
2004--2005 \quad Researcher, Rigshospitalet, Denmark \\
1996--2004 \quad Researcher/Postdoctoral Researcher, Technical University of Denmark, Denmark

\section*{Education}
2001 \quad PhD, Technical University of Denmark, Denmark \\
1996 \quad Master of Engineering, Technical University of Denmark, Denmark \\
1993 \quad Degree, Aarhus University School of Engineering, Denmark

\section*{Career breaks}
Not relevant

\section*{Research statement}
My research focuses on the intersection of semantic technologies, linked open data, and scientific knowledge organization. I have pioneered the development of tools and workflows that leverage Wikidata and Scholia to visualize and analyze complex scientific networks and bibliographic data. My work transitions from neuroimaging database development to the creation of semantic representations for large-scale knowledge bases. I aim to enhance the reproducibility and discoverability of science by building robust, machine-readable infrastructures that support living reviews and automated research workflows.

\section*{Personal context}
I hold a long-standing academic position at the Technical University of Denmark, where my role involves a balance of high-level research development, software engineering for scientific infrastructures, and academic responsibilities. My research trajectory has evolved from neuroimaging data analysis to the broader, multidisciplinary field of open science and semantic web technologies, requiring significant time for both software maintenance and community-driven development.

\section*{Grants and awards}
2007 \quad Jorck's Prize, [Insert Funder/Source]

\section*{Supervision, teaching and research leadership}
I have significant experience in academic mentoring and research leadership, providing guidance to various researchers. 
\begin{itemize}
    \item \textbf{Supervision:} Supervised students including Laura Rieger; provided academic advisement to Lars Kai Hansen and Jan Larsen.
    \item \textbf{Leadership:} Experience in leading technical projects related to the development of semantic software and research tools.
\end{itemize}

\section*{Collaborations and teamwork}
I maintain an extensive international network of collaborators across academia and technical communities. Key collaborators include Adélie Ranville, Andra Waagmeester, Daniel Mietchen, Egon Willighagen, Claus Svarer, Lars Kai Hansen, Dario Taraborelli, Najko Jahn, Gitte Moos Knudsen, Rémy Gerbet, Dariusz Jemielniak, and Arthur Perret. My work often involves large-scale scientific teamwork, particularly in the context of open-source software development and linked data ecosystems.

\section*{Contributions to the research community}
I am a dedicated contributor to the infrastructure of open science. My contributions include the development of Scholia for visualizing scholarly data and the advancement of WikiCite to improve bibliographic citations in the sum of all human knowledge. I have developed specialized software for creating and analyzing semantic representations and have worked extensively on making scientific data (such as the EU NanoSafety Cluster) more transparent through Linked Open Data visualization.

\section*{Contributions to the wider society}
I contribute to society by promoting the democratization of knowledge through the use of Wikimedia-integrated workflows. My work on "Games for the lexemes in Wikidata" and the development of tools for sustainability transitions demonstrates a commitment to making complex data accessible and engaging for the general public, thereby supporting science communication and community engagement.

\end{document}